\section{实验}
\label{sec:experiments}

\subsection{实验设置}

\textbf{评测基准。}采用ClawEval的纯文本子集，共195个任务，按表\ref{tab:buckets}的九个能力域分组。评分采用Pass$^3$协议（三次独立运行全部通过），逐任务二元判定、逐桶通过率。

\textbf{训练协议。}训练\textbf{按桶顺序}进行：模型先在一个能力域上训练，再切换到下一个，模拟线上无法一次集齐全分布的现实约束。同一桶内任务随机打乱。基座模型为Qwen3.5-9B，采用完全异步的策略分离部署（推理与训练分卡），BF16全栈。

\textbf{算法配置。}策略采用GRPO，每组8条查询轨迹。裁判模型为独立部署的冻结外部LLM。底层RL框架verl 0.8.0不做任何改动，所有持续学习逻辑通过损失注入接口与缓冲区钩子挂载。

\subsection{抗遗忘评测协议}

全部能力域顺序训完后\textbf{一次统一评测}（训练过程中不做中间评测，省算力）：

\begin{enumerate}
  \item 对终态策略$\pi_K$跑一次全量评测。
  \item 评测结果按同样的九个桶分组，组间序与训练序对齐，记录每桶原始分数$S_i = S(\pi_K, b_i)$。
  \item 施加按训练先后\textbf{衰减的权重}$w_i$：越早训练的桶遗忘越重、权重越大：
  \begin{equation}
    \text{CL-Score} = \sum_{i=1}^{K} w_i\,S_i,\qquad w_i \propto g(i),\;\;g\text{递减}.
  \end{equation}
  衰减形式（线性/指数）作为超参由实验确定。
\end{enumerate}

评测只跑一次，之后切换权重方案只是算术运算，无需重新推理。核心信号是含持续学习机制的模型在\textbf{更早训练的桶上}相对基线表现出\textbf{更小的分数下降}。

\subsection{对照与消融}

\textbf{核心对照}：含持续学习机制（九桶回放+复合损失）\emph{vs.} 无持续学习机制的同规模基线（无回放、$\lambda_3=0$，相同训练序和评测）。

\textbf{消融轴}：
\begin{itemize}
  \item 缓冲区结构：九桶分区回放 \emph{vs.} CLEAR式单扁平缓冲区+蓄水池采样~\cite{A2}。
  \item 缓冲区容量：不同总容量下的抗遗忘表现。
  \item U形加权：$\gamma=\delta=1$（均权，默认）\emph{vs.} $\gamma=\delta<1$（U形激活）。
\end{itemize}

\subsection{实验结果}

\todo{新reward+CL算法16卡正式实验就绪后填入。预期产出：逐桶原始分数；CL-Score对照表（CL vs. baseline）；reward/loss/组内std训练曲线；消融柱状图。}
